\documentclass[10pt,a4paper]{amsart}
\usepackage[margin=19mm]{geometry}
\usepackage{amsmath,amssymb,mathtools}
\usepackage[scr=rsfs,cal=euler]{mathalfa}
\usepackage{xfrac}
\usepackage[unicode,hidelinks,bookmarks=false]{hyperref}
\newcommand{\ie}{i.e.\ }
\newcommand{\cf}{cf.\ }
\newcommand{\resp}{respectively\ }
\newcommand{\Subsection}{\subsection}
\providecommand{\nobreakdash}{\nobreak\hbox{-}}
\setlength{\emergencystretch}{2em}
\multlinegap0pt
\usepackage{bm}
\usepackage{bbm}
\usepackage{dsfont}
\usepackage{xspace}
\usepackage{cancel}
\usepackage{mathtools}
\usepackage{amsthm}
\makeatletter
\@ifundefined{Lemma}{\newtheorem{Lemma}{Lemma}}{}
\@ifundefined{Proposition}{\newtheorem{Proposition}[Lemma]{Proposition}}{}
\@ifundefined{Corollary}{\newtheorem{Corollary}[Lemma]{Corollary}}{}
\makeatother
\title[Some formulae relating modular representations of elementary]{Some formulae relating modular representations of elementary abelian $p$-groups}
\author{Jonathan Elmer, Kazal Kadr}
\date{Source: arXiv:2510.07939v1 (CC BY 4.0)}
\begin{document}
\maketitle

\noindent\emph{Source and licence.} Some formulae relating modular representations of elementary abelian $p$-groups. Version arXiv:2510.07939v1 (CC BY 4.0), distributed under the Creative Commons Attribution 4.0 International licence, as recorded in the arXiv licence metadata for this exact version and shown on the abstract page https://arxiv.org/abs/2510.07939v1. Full rights notice: licenses/arxiv-2510.07939-CC-BY-4.0.txt.\par\smallskip

\noindent\textit{Editorial scope.} Counted unit: the Lemma whose source label is \texttt{notenough} in the article (source0.tex lines 601--602, source label \texttt{notenough}), delivered together with the article's own complete original proof of it (source0.tex lines 604--638, one \texttt{proof} environment of 35 lines). No mathematics is written, repaired or re-derived: statement and proof are the article's own text line for line. Required context retained uncounted: newtensorformula (lines 164--167); wildon (lines 304--306); dualaction (lines 322--323); morecombs (lines 640--642). Every definition, display and label of those lines is retained unchanged. Omitted from this package and not redistributed: 1--163, 168--303, 307--321, 324--600, 603--603, 639--639, 643--687. Nothing inside a retained excerpt is omitted or altered: every retained body is delivered exactly as the article's lines are written, so no omission receipt of any kind accompanies this package. Citations: the retained excerpts carry no citation command at all, so no bibliography entry of the article is transcribed and no reference is redistributed. Reference adaptation: none was needed and none was made; every reference inside the delivered unit resolves inside it, so no unresolved reference or citation is produced. Printed slips kept literal: no printed word, symbol, hypothesis or number of the counted statement or of its proof was altered. Layout and class: the article's own class is \texttt{amsart} with packages including latexsym, amssymb, amsmath, amsthm, amsfonts, tikz, mathtools; the wrapper is the lane's pinned \texttt{amsart} editorial wrapper at 10pt on A4, which restates the theorem-like environments the retained text needs and adds the article's own macro definitions that the retained text needs (none), with extra wrapper packages (bm, bbm, dsfont, xspace, cancel, mathtools). The delivered numbering is the unit's own. Assets: no retained span contains a figure environment, a graphics-inclusion command, a PDF-page inclusion, a table or a TikZ picture; no image file is copied into this package, the package declares no asset dependency and no image file is redistributed with it. Licence: version arXiv:2510.07939v1 of this article is distributed under the Creative Commons Attribution 4.0 International licence, as recorded in the arXiv licence metadata for this exact version and shown on the abstract page https://arxiv.org/abs/2510.07939v1. A full rights notice is delivered at licenses/arxiv-2510.07939-CC-BY-4.0.txt. Characters: every retained excerpt is pure ASCII, so the delivered engine, XeLaTeX with its default Latin Modern text font, reports no missing character.
\begin{Corollary}\label{newtensorformula}
Let $i<p$ and $j<q$. Assume $i+j \leq q$. Write $j = rp+ j'$ and assume that $i+j' \leq p$. Then we have
\begin{equation}\label{newtensorbyVj} V_i \otimes V_j \cong_M \bigoplus_{l=1}^{\min(i,j')} V_{j+i-(2l-1)}\end{equation}
\end{Corollary}

\begin{Proposition}\label{wildon}
Let $d,i>0$. Then  $S^d(V_{i+1}) \cong S^i(V_{d+1}^*)^*$ for all $d,i < q$, and $S^d(V_{i+1})  \cong \Lambda^d(V_{i+d})$ if in addition $d+i \leq q$.
\end{Proposition}

\begin{equation}\label{dualaction} \alpha \cdot x_i  = \sum_{j=0}^{m-i}\binom{i+j}{i}(-\alpha)^j x_{i+j}.
\end{equation}

\begin{Lemma}\label{morecombs}
For all $i,s,t$ we have  $\sum_{j+k=i} \binom{t}{j}\binom{s}{k} =  \binom{s+t}{i}$.
\end{Lemma}

\begin{Lemma}\label{notenough} Suppose that $p>2$. Then we have $V_{p+1} \otimes V_{p+1} \cong_M V_{2p+1}$.
\end{Lemma}

\begin{proof}
We first note that since $p>2$ we have  $V_{p+1} \otimes V_{p+1} \cong S^2(V_{p+1}) \oplus \Lambda^2(V_{p+1})$. Now Proposition \ref{wildon} implies that $ \Lambda^2(V_{p+1}) \cong S^2(V_p)$. The latter is a direct summand of $V_p \otimes V_p$ which is projective relative to $M$. Therefore 
$V_{p+1} \otimes V_{p+1} \cong_M S^2(V_{p+1})$ and it is enough to show that $S^2(V_{p+1}) \cong_M V_{2p+1}$.

We claim that  $S^2(V_{p+1}) \cong S^2(V_{p-1}) \oplus V_{2p+1}$. We will exhibit a submodule of $V_{p+1}$ isomorphic to $V_{p-1}$. Let $x_0, x_1, \ldots, x_{p}$ be the basis of $V_{p+1}$ with respect to which the action of $\alpha \in G$ is given by the usual formula \eqref{dualaction}. For $i=0, \ldots p-2$ we set $y_i = (i+1)x_{i+1}$. Then we have
\begin{align*} \alpha \cdot y_i = (i+1) \alpha \cdot x_{i+1} &= (i+1) \sum_{k=0}^{p-1-i} \binom{i+k+1}{k}(-\alpha)^k x_{i+1+k}. \end{align*}
Note that the coefficient of $x_{p}$ on the right hand side is $\binom{p}{i+1} = 0 \mod p$, so the sum only needs to run to $p-2-i$. Thus 
\begin{align*}\alpha \cdot y_i &=  (i+1) \sum_{k=0}^{p-2-i} \binom{i+1+k}{k}(-\alpha)^k \frac{1}{i+1+k}y_{i+k} \\
&=  \sum_{k=0}^{p-2-i} \binom{i+k}{k}(-\alpha)^k y_{i+k} \end{align*}
as required.

It follows that $S^2(V_{p+1})$ contains a submodule $Y$ isomorphic to $S^2(V_{p-1})$, with basis $\{x_ix_j: i,j = 1, \ldots, p-1\}$.
Now we define, for each $i=0, \ldots, 2p$
\[z_i = (-1)^i \sum_{j+k=i}x_jx_k.\]

We claim that $z_0, z_1, \ldots, z_{2p}$ span a submodule of $S^2(V_{p+1})$ isomorphic to $V_{2p+1}$. 
We have, for all $i=0, \ldots, 2p$

\begin{align*}
\alpha \cdot z_i &= (-1)^i \sum_{j+k=i}(\alpha \cdot x_j)(\alpha \cdot x_k)\\
&= (-1)^i \sum_{j+k=i}\left(\sum_{t=0}^p \binom{t}{j} (-\alpha)^{t-j} x_t \right) \left(\sum_{s=0}^p \binom{s}{k} (-\alpha)^{s-k} x_s \right)\\
&= (-1)^i \sum_{j+k=i} \sum_{r=0}^{2p} \sum_{s+t=r}\left(\binom{t}{j} (-\alpha)^{t-j} x_t \right) \left(\binom{s}{k} (-\alpha)^{s-k} x_s \right)\\
&= (-1)^i \sum_{j+k=i} \sum_{r=0}^{2p} (-\alpha)^{r-i} \sum_{s+t=r} \binom{t}{j}\binom{s}{k}  x_s x_t\\
&= (-1)^i  \sum_{r=0}^{2p} (-\alpha)^{r-i} \sum_{s+t=r}\left(  \sum_{j+k=i} \binom{t}{j}\binom{s}{k} \right)  x_s x_t\\
\intertext{by Lemma \ref{morecombs} below}
&= (-1)^i  \sum_{r=0}^{2p} (-\alpha)^{r-i} \sum_{s+t=r} \binom{r}{i} x_s x_t\\
&=\sum_{r=0}^{2p} (-\alpha)^{r-i} \binom{r}{i} z_r
\end{align*}
as required.

This proves our claim. Now to finish the proof, observe that $Z \cap Y = \{0\}$, since $z_i$ contains a unique term $2x_0x_i$ for $i \leq p$ or $2x_{i-p}x_p$ for $i>p$, none of which are contained in any term in $Y$ (again we need $p>2$ here). Thus, since $\dim(Z)+\dim(Y) = 2p+1+\frac12 p(p-1) = \frac12 (p+1)(p+2) = \dim(S^2(V_{p+1}))$ we have $S^2(V_{p+1}) = Y \oplus Z \cong S^2(V_{p-1}) \oplus V_{2p+1}$.

Finally we claim that $S^2(V_{p-1}) \cong_M 0$. To see this, note that $S^2(V_{p-1})$ is a direct summand of $V_{p-1} \otimes V_{p-1}$. By Corollary \ref{newtensorformula} we have $V_{p-1} \otimes V_{p-1} \cong_M V_1$. But as $\dim( S^2(V_{p-1}) ) = \frac12 p(p-1) = 0 \mod p$, we cannot have 
$S^2(V_{p-1}) \cong_M V_1$ so we must have  $S^2(V_{p-1}) \cong_M 0$ as required. Alternatively we can use the isomorphism $S^2(V_{p-1}) \cong \Lambda^2(V_p)$, and the latter is a summand of $V_p \otimes V_p$. This completes the proof.
\end{proof}

\end{document}
